\section{Xingmian (星眠) y Dokie: dos líneas de producto}

Ya se explicó la decisión emprendedora de apostarlo todo (all in); ahora hay que ver en qué productos concretos se materializa esa decisión. Xingmian y Dokie representan dos rutas: el contenido de compañía y el Agent de productividad.

Muyan Zhiyu (沐言智语) tiene hoy dos líneas principales: el juego otome con IA Xingmian y el producto Agent Dokie. El primero continúa el interés de largo plazo de Zhang Yueguang (张月光) por Live2D, la compañía y el contenido dirigido al público femenino; el segundo es la dirección de productividad en la que decidió apostarlo todo (all in) tras un año de exploración.

\begin{figure}[H]
\centering
\includegraphics[width=0.94\textwidth]{xingmian-dokie-two-bets.png}
\caption{Las dos líneas de producto de Muyan Zhiyu: Xingmian y Dokie.}
\end{figure}

\begin{knowledgebox}{Lectura de la figura: uno hace contenido, el otro amplía los límites de capacidad}
Xingmian pone a prueba el contenido y la compañía con IA; lo que importa ahí es el contenido, los personajes, la frecuencia y la interacción. Dokie entra por las PPT y la expresión creativa, y su objetivo es que el usuario supere los límites de su capacidad. Ninguno de los dos es una simple demostración técnica: ambos buscan una puerta de entrada con la que el usuario interactúe, a la que recuerde y por la que pague a largo plazo.
\end{knowledgebox}

\subsection{Por qué los juegos con IA son un espacio estrecho}

Zhang Yueguang considera que el espacio viable de los juegos con IA es más estrecho de lo que mucha gente cree. La generación de video en tiempo real, los juegos cinematográficos de interacción infinita y las tramas dinámicas en mundo abierto todavía son difíciles de llevar a la práctica de forma estable con la tecnología actual. Donde la IA aporta valor con más facilidad es en la conversación de compañía y en la relación con los personajes; por eso confía más en los juegos otome con IA que en los juegos con IA en general.

\begin{warningbox}{Un juego con IA no es cualquier juego más IA}
En esencia, los juegos siguen siendo una industria de contenido. La IA puede cambiar ciertas experiencias de contenido, pero no resuelve por sí sola la trama, los costos, la estabilidad ni la retención de usuarios. Los escenarios viables suelen ser más estrechos que las grandes visiones.
\end{warningbox}

\subsection{Por qué Dokie entra por las PPT}

Dokie entra por las PPT, pero Zhang Yueguang insiste en que no será para siempre una herramienta de PPT. Las PPT son una puerta de entrada a la creación y la expresión; el objetivo real es «un amigo de IA que me ayude a superar los límites de mi capacidad». Esto explica por qué está dispuesto a apostarlo todo (all in) por el Agent y no a poner todo el capital en aquellas exploraciones pequeñas de antes.

\subsection{Resumen de la sección}

Xingmian y Dokie representan dos formas de la población de IA: el individuo de IA en la compañía y el contenido, y el individuo de IA en la productividad y la creación. El primero tiene que controlar el costo del contenido y los límites de la compañía; el segundo tiene que demostrar que de verdad se superan los límites de capacidad.
